\chapter{Distillation: 큰 경험을 작은 상태로 옮기기}
\label{bg:chapter-distillation}

\section{교사와 학생}

\concept{distillation}{지식 증류}는 \concept{teacher-model}{교사 모델}이 제공한 output
distribution, rationale, feature, selected answer를 \concept{student-model}{학생 모델}이
모방하도록 학습하는 방식이다. 교사가 더 크거나 더 많은 context를 사용할 수 있고, 학생은
작거나 latency가 낮을 수 있다. Sleep context에서는 긴 episodic history를 짧은 memory,
adapter, 작은 model로 압축하는 operator로 이해할 수 있다.

Hard target은 교사의 최고 확률 token 하나만 사용한다. Soft target은 전체 분포를 사용해
두 번째·세 번째 선택지에 대한 정보도 전달한다. \concept{temperature}{온도} $T$를 적용하면

\begin{equation}
  p_i^{(T)}=\frac{\exp(z_i/T)}{\sum_j\exp(z_j/T)}.
\end{equation}

$T>1$에서 분포가 평평해져 relative preference가 더 드러난다. 학생 분포 $q$와 교사 분포
$p$ 사이 \concept{kl-divergence}{KL 발산}을 줄일 수 있다.

\begin{equation}
  D_{\mathrm{KL}}(p\Vert q)=\sum_i p_i\log\frac{p_i}{q_i}.
  \label{eq:bg-kl}
\end{equation}

KL은 대칭 거리가 아니다. 어떤 방향을 쓰는지에 따라 mode covering과 mode seeking 특성이
달라진다. 실제 distillation loss는 hard label loss와 soft loss를 섞는다.

\begin{equation}
  L=\alpha L_{\mathrm{hard}}+(1-\alpha)T^2
  D_{\mathrm{KL}}(p_T\Vert q_T).
  \label{eq:bg-distill-loss}
\end{equation}

\section{Memory consolidation으로서 distillation}

긴 history를 매번 prompt에 넣는 대신 teacher가 충분한 context와 retrieval을 사용해 reusable
summary나 strategy를 만든다. destination에 따라 세 종류로 나뉜다.

\begin{enumerate}
  \item \textbf{External distillation}: episode를 concise text, vector prototype, graph fact로 바꾼다.
  \item \textbf{Adapter distillation}: teacher의 behavior를 LoRA나 adapter에 학습한다.
  \item \textbf{Model distillation}: 여러 cycle의 memory를 base model weights로 통합한다.
\end{enumerate}

첫 번째는 provenance와 수정이 쉽고 inference 때 retrieval cost가 남는다. 두 번째는 repeated
query의 retrieval cost를 줄일 수 있지만 base model compatibility와 adapter routing이 필요하다.
세 번째는 access latency가 가장 낮을 수 있으나 interference, deletion, rollback cost가 가장
크다.

\section{무엇이 압축 과정에서 사라지는가}

Distillation은 loss가 중요하게 보는 behavior만 보존한다. rare exception, uncertainty,
source wording, provenance는 teacher output에 드러나지 않으면 사라진다. 외부 summary가
``사용자는 항상 아침 운동을 선호한다''고 일반화했지만 원 episode가 ``이번 주만''이었다면
identity drift가 발생한다.

세 가지 검증을 분리한다.

\begin{itemize}
  \item \textbf{fidelity}: teacher와 student가 같은 query distribution에서 얼마나 같은가?
  \item \textbf{utility}: teacher와 달라도 downstream task를 실제로 더 잘 푸는가?
  \item \textbf{provenance preservation}: output의 각 claim을 source episode로 되돌릴 수 있는가?
\end{itemize}

Fidelity가 너무 높으면 teacher 오류까지 복사한다. utility만 최적화하면 근거 없는 shortcut을
배울 수 있다. provenance는 loss 한 항으로 쉽게 보장되지 않으므로 external metadata와
release gate가 필요하다.

\section{Recursive distillation과 model collapse}

Sleep cycle $k$의 output이 cycle $k+1$의 training data가 되면 recursive depth가 증가한다.
실제 data anchor가 줄고 생성 data만 반복되면 distribution tail이 사라질 수 있다
\parencite{shumailov2024collapse}. 반면 real data를 계속 누적하는 조건에서는 collapse가
필연적이지 않을 수 있다. 따라서 synthetic 여부의 이분법 대신 각 item에
origin, teacher version, recursive depth, real-evidence links를 기록한다.

\begin{table}[htbp]
\centering
\caption{Sleep distillation destination의 차이}
\label{tab:bg-distillation-destinations}
\begin{tabularx}{\textwidth}{L{0.19\textwidth}YYYY}
\toprule
Destination & update 비용 & wake 비용 & 수정/삭제 & 대표 위험 \\
\midrule
text summary & 낮음--중간 & retrieval/context & 쉬움 & over-compression \\
vector prototype & 중간 & ANN search & 중간 & embedding drift \\
temporal graph & 중간--높음 & traversal + context & 비교적 쉬움 & conflict/edge explosion \\
LoRA/adapter & training 필요 & routing + matmul & version rollback 가능 & conflict/base churn \\
full weights & 가장 높음 & 낮을 수 있음 & 가장 어려움 & interference/unlearning \\
\bottomrule
\end{tabularx}
\end{table}

\begin{keypoint}
Distillation은 기억을 ``옮기는'' 마법이 아니라 task distribution이 정한 정보만 보존하는
lossy compression이다. Sleep system은 압축률뿐 아니라 삭제 가능한 provenance와 later
utility를 함께 최적화해야 한다.
\end{keypoint}
